% !TEX root = ../memoria.tex
\chapter{Análisis y Diseño}
\label{cap:capitulo3}

Este capítulo describe las decisiones de análisis y diseño que guían la construcción del
gemelo digital de cobertura 5G. Se parte de la elicitación de requisitos funcionales y no
funcionales, se define la arquitectura en tres capas del sistema, se detalla el modelo de
datos JSON que conecta la simulación con la visualización y se expone el diseño de la API
REST y de la interfaz de usuario. El capítulo cierra con la justificación razonada del
\textit{stack} tecnológico y del modelo de propagación seleccionado.

% ============================================================
\section{Requisitos del Sistema}
\label{sec:requisitos}
% ============================================================

La especificación de requisitos se ha elaborado siguiendo un enfoque incremental: los
requisitos iniciales se extrajeron de las necesidades del caso de uso táctico descrito en
el Capítulo~\ref{cap:capitulo2} y se refinaron durante las iteraciones de desarrollo. Se
separan en requisitos funcionales (qué debe hacer el sistema) y requisitos no funcionales
(con qué restricciones de calidad debe hacerlo).

\subsection{Requisitos Funcionales}
\label{subsec:req_func}

Los requisitos funcionales definen las capacidades que el sistema debe proporcionar al
usuario final ---el operador táctico que planifica el despliegue de una burbuja 5G--- y al
ingeniero de telecomunicaciones que analiza los resultados de la simulación.

\begin{table}[H]
\centering
\caption{Requisitos funcionales del sistema.}
\label{tab:req_func}
\begin{tabular}{|c|p{10cm}|c|}
\hline
\textbf{ID} & \textbf{Descripción} & \textbf{Prioridad} \\
\hline
RF1 & El sistema debe simular la cobertura 5G FR1 (banda n78, 3,5~GHz)
      aplicando el modelo de propagación Longley-Rice ITM sobre datos de
      elevación SRTM de 30~m de resolución. & Alta \\
\hline
RF2 & Los resultados de cada simulación deben exportarse en formato JSON
      con un identificador UUID~v4 único, incluyendo parámetros de
      configuración, metadatos y los puntos de cobertura con coordenadas
      y valores RSRP. & Alta \\
\hline
RF3 & El servidor debe exponer una API REST que permita consultar el
      listado de simulaciones disponibles y obtener el detalle completo
      de cualquiera de ellas mediante su UUID. & Alta \\
\hline
RF4 & La aplicación web debe visualizar los puntos de cobertura de una
      simulación como un \textit{heatmap} interactivo superpuesto sobre
      un mapa cartográfico, con colores representativos de los niveles
      de RSRP. & Alta \\
\hline
RF5 & El usuario debe poder filtrar de forma interactiva los puntos de
      cobertura por umbral de RSRP mediante un control deslizante
      (\textit{slider}), actualizándose el mapa en tiempo real. & Media \\
\hline
\end{tabular}
\end{table}

\subsection{Requisitos No Funcionales}
\label{subsec:req_nofunc}

Los requisitos no funcionales establecen las restricciones de calidad, rendimiento y
compatibilidad que el sistema debe satisfacer para ser aceptable en un entorno de uso
táctico real, donde la rapidez y la fiabilidad son críticas.

\begin{table}[H]
\centering
\caption{Requisitos no funcionales del sistema.}
\label{tab:req_nofunc}
\begin{tabular}{|c|p{10cm}|c|}
\hline
\textbf{ID} & \textbf{Descripción} & \textbf{Categoría} \\
\hline
RNF1 & La interfaz de usuario debe ser responsiva, adaptándose
       correctamente a resoluciones de escritorio ($\geq 1280$~px) y
       tableta ($\geq 768$~px). & Usabilidad \\
\hline
RNF2 & El tiempo de carga y renderizado del \textit{heatmap} para una
       simulación de 3600~puntos no debe superar los 3~segundos en un
       equipo de gama media con conexión local. & Rendimiento \\
\hline
RNF3 & El backend debe ser capaz de servir entre 5 y 10 ficheros de
       simulación simultáneamente sin degradación perceptible del
       rendimiento. & Escalabilidad \\
\hline
RNF4 & La aplicación web debe funcionar correctamente en las dos últimas
       versiones estables de Google Chrome, Mozilla Firefox y Microsoft
       Edge. & Compatibilidad \\
\hline
RNF5 & La cobertura de tests unitarios y de integración del backend debe
       ser superior al 60~\% del código fuente. & Mantenibilidad \\
\hline
\end{tabular}
\end{table}

% ============================================================
\section{Arquitectura del Sistema}
\label{sec:arquitectura}
% ============================================================

La arquitectura del sistema sigue un patrón de tres capas desacopladas, cada una con una
responsabilidad bien definida: simulación, servicio de datos y presentación. Este diseño
permite que cada capa evolucione de forma independiente y facilita la sustitución de
componentes sin impacto en el resto del pipeline.

\subsection{Visión General de la Arquitectura}
\label{subsec:arq_general}

El flujo de datos del gemelo digital es unidireccional y consta de cinco etapas, como se
muestra a continuación:

\begin{enumerate}
  \item \textbf{Simulación} (MATLAB, ejecución \textit{offline}): el motor de simulación
    calcula la cobertura RSRP sobre una cuadrícula de $60 \times 60$~puntos utilizando el
    modelo Longley-Rice ITM con datos de elevación SRTM. El resultado se serializa en un
    fichero JSON nombrado con un UUID~v4.
  \item \textbf{Almacenamiento}: los ficheros JSON se depositan en un directorio compartido
    (\lstinline{matlab/output/}) que actúa como almacén persistente del sistema. No se
    emplea base de datos relacional: el fichero JSON \emph{es} la fuente de verdad.
  \item \textbf{Servicio} (Spring Boot): la API REST lee los ficheros JSON del directorio
    configurado, los deserializa con Jackson y los expone mediante \textit{endpoints}
    RESTful documentados con OpenAPI.
  \item \textbf{Transporte}: las peticiones HTTP entre frontend y backend se realizan sobre
    \lstinline{localhost} en desarrollo (puertos 4200 y 8080 respectivamente), con CORS
    habilitado explícitamente.
  \item \textbf{Presentación} (Angular + Leaflet): la aplicación web consume la API,
    transforma los datos al modelo interno y renderiza el \textit{heatmap} interactivo
    sobre un mapa OpenStreetMap.
\end{enumerate}

Este diseño desacoplado tiene una ventaja importante para el caso de uso táctico: la capa
de simulación puede ejecutarse en un portátil con MATLAB sin conexión a internet, generar
los JSON con los resultados y, posteriormente, cargarlos en un servidor Spring Boot ligero
que sirva la visualización a cualquier dispositivo conectado a la red local de la burbuja
5G. No se requiere conectividad externa en ningún punto del pipeline.

\subsection{Capa de Simulación: Motor MATLAB}
\label{subsec:arq_matlab}

El motor de simulación es un script MATLAB (\lstinline{simulacion_cobertura_5g.m})
estructurado en nueve bloques ejecutables de forma independiente. Su responsabilidad es
triple: (1)~configurar los parámetros de la estación base y la cuadrícula de receptores;
(2)~seleccionar automáticamente el mejor modelo de propagación disponible en función de los
\textit{toolboxes} instalados; y (3)~exportar los resultados en formato JSON compatible con
el backend.

La selección de modelo sigue una lógica de degradación controlada. Si el
\textit{Communications Toolbox} está disponible, se utiliza el modelo Longley-Rice ITM, que
incorpora el perfil de elevación real del terreno obtenido de datos SRTM a 30~m de
resolución. Si el \textit{toolbox} no está disponible, el script recurre al modelo
COST-231 Hata como \textit{fallback}, registrando una advertencia en los metadatos del JSON
de salida. En ambos casos se calcula también la pérdida en espacio libre (FSPL) como cota
superior teórica de validación.

La integración de datos de elevación se realiza mediante el \textit{Mapping Toolbox}:
las funciones \lstinline{readgeoraster} y \lstinline{geointerp} cargan el \textit{tile}
SRTM (GeoTIFF) e interpolan la altitud en cada punto de la cuadrícula. Si el fichero
SRTM no está disponible, el sistema asume terreno plano y continúa con la simulación,
marcando el campo \lstinline{terrain_data} del JSON como \lstinline{"flat_assumed"}.

\subsection{Capa de Servidor: API REST Spring Boot}
\label{subsec:arq_backend}

El backend está implementado con Spring Boot 3.2.4 sobre Java~17, utilizando Maven como
gestor de dependencias. Su función es exponer los resultados de las simulaciones MATLAB a
través de una API REST \textit{stateless}, sin persistencia en base de datos: los ficheros
JSON del directorio de salida de MATLAB son la única fuente de datos.

La estructura de paquetes sigue las convenciones estándar de Spring Boot:

\begin{itemize}
  \item \lstinline{controller/}: clase \lstinline{SimulationController} con los
    \textit{endpoints} REST anotados con \lstinline{@RestController}.
  \item \lstinline{service/}: clase \lstinline{SimulationService} que encapsula la lógica
    de lectura y deserialización de ficheros JSON.
  \item \lstinline{model/}: \textit{records} Java (\lstinline{SimulationSummaryDto},
    etc.) que representan los DTOs de respuesta.
  \item \lstinline{exception/}: clase \lstinline{SimulationNotFoundException} anotada con
    \lstinline{@ResponseStatus(HttpStatus.NOT_FOUND)} y manejador global de excepciones.
  \item \lstinline{config/}: clase \lstinline{CorsConfig} que habilita las peticiones
    \textit{cross-origin} desde \lstinline{localhost:4200}.
\end{itemize}

El diseño RESTful es completamente \textit{stateless}: cada petición contiene toda la
información necesaria para ser procesada, no se mantiene estado de sesión en el servidor y
todos los \textit{endpoints} son idempotentes (operaciones \lstinline{GET} exclusivamente).
La documentación interactiva de la API se genera automáticamente mediante
\lstinline{springdoc-openapi 2.3.0}, accesible en \lstinline{/swagger-ui.html}.

\subsection{Capa de Presentación: Aplicación Angular}
\label{subsec:arq_frontend}

La interfaz de usuario está desarrollada con Angular~17+ utilizando componentes
\textit{standalone} (sin \lstinline{NgModule}), lo que simplifica la estructura del proyecto
y mejora el \textit{tree-shaking} del compilador. La aplicación consume la API REST del
backend y presenta los datos de cobertura como un \textit{heatmap} interactivo renderizado
sobre un mapa OpenStreetMap mediante las bibliotecas Leaflet y \lstinline{leaflet.heat}.

La gestión de estado se basa en \textit{signals} de Angular, un mecanismo reactivo
introducido en Angular~16 que permite propagar cambios de forma eficiente sin necesidad de
bibliotecas externas como NgRx o RxJS para el estado local. El servicio
\lstinline{SimulationService} expone \textit{signals} de solo lectura para el estado de
carga (\lstinline{loading}), los errores (\lstinline{error}) y la lista de simulaciones
disponibles.

La comunicación con el backend se realiza mediante \lstinline{HttpClient} de Angular,
inyectado mediante \lstinline{provideHttpClient()} en la configuración de arranque de la
aplicación. Los DTOs del backend (\lstinline{ApiSimulation}, \lstinline{ApiCoveragePoint})
se transforman al modelo interno de la aplicación mediante la función
\lstinline{apiToSimulation()}, desacoplando así la representación del servidor de la
representación de la vista.

% ============================================================
\section{Diseño del Modelo de Datos}
\label{sec:datos}
% ============================================================

El fichero JSON constituye el contrato de datos entre la capa de simulación y el resto del
sistema. Su estructura ha sido diseñada para ser autocontenida: cada fichero incluye tanto
los parámetros de configuración de la simulación como los resultados, de forma que puede
interpretarse sin información adicional externa.

\subsection{Esquema JSON de Simulación}
\label{subsec:esquema_json}

Cada simulación se almacena en un fichero JSON independiente, nombrado con su UUID~v4
(por ejemplo, \lstinline{b83baf94-8db6-4512-a34b-b831e3f6267a.json}). El esquema del
fichero se muestra en el Listado~\ref{cod:json_schema}.

\begin{lstlisting}[caption={Estructura del fichero JSON de simulación.},
  label={cod:json_schema}, basicstyle={\footnotesize\ttfamily}]
{
  "id":        "b83baf94-8db6-4512-a34b-b831e3f6267a",
  "name":      "URJC Fuenlabrada - Longley-Rice ITM",
  "createdAt": "2026-04-10T14:32:00Z",
  "parameters": {
    "frequency_hz":    3500000000,
    "tx_power_dbm":    40,
    "antenna_height_m": 30,
    "rx_height_m":     1.5,
    "tx_latitude":     40.2897,
    "tx_longitude":    -3.8244,
    "area_km":         3,
    "espaciado_grid_m": 50
  },
  "metadata": {
    "frequency_ghz":    3.5,
    "tx_power_dbm":     40,
    "antenna_height_m": 30,
    "tx_location":      { "latitude": 40.2897, "longitude": -3.8244 },
    "simulation_date":  "2026-04-10",
    "model":            "longley-rice",
    "terrain_data":     "srtm_30m"
  },
  "coveragePoints": [
    { "lat": 40.2762, "lng": -3.8379, "rsrp": -97.3 },
    { "lat": 40.2762, "lng": -3.8374, "rsrp": -95.1 },
    ...
  ]
}
\end{lstlisting}

Los campos principales son:

\begin{itemize}
  \item \lstinline{id}: identificador UUID~v4 generado por MATLAB en el momento de la
    exportación. Sirve como clave primaria en el backend y como nombre del fichero.
  \item \lstinline{parameters}: parámetros de configuración de la simulación (frecuencia,
    potencia, alturas de antena, coordenadas del emplazamiento y geometría de la cuadrícula).
  \item \lstinline{metadata}: información descriptiva para el frontend (modelo utilizado,
    tipo de datos de terreno, fecha de ejecución).
  \item \lstinline{coveragePoints}: array de objetos con las coordenadas geográficas
    (\lstinline{lat}, \lstinline{lng}) y el valor RSRP en dBm de cada punto de la
    cuadrícula. Una simulación estándar contiene 3600~puntos.
\end{itemize}

\subsection{Validación de Datos}
\label{subsec:validacion_datos}

La integridad de los datos se garantiza mediante validaciones en dos niveles del sistema.

En la \textbf{capa de simulación}, MATLAB aplica restricciones sobre los valores generados
antes de la exportación. Los valores de RSRP se recortan al rango $[-140, -40]$~dBm, que
corresponde a los límites definidos en 3GPP TS~38.215 \cite{3gpp_ts38215_2022} para
medidas de potencia de señal de referencia. La frecuencia de simulación se restringe al
rango de la banda n78 (3,4--3,8~GHz), verificándose que el valor configurado en
\lstinline{fc_Hz} caiga dentro de este intervalo.

En la \textbf{capa de servidor}, el backend Spring Boot aplica validaciones adicionales
sobre las peticiones entrantes:

\begin{itemize}
  \item \textbf{Validación de UUID}: el parámetro de ruta \texttt{\{uuid\}} se valida
    contra una expresión regular que verifica el formato UUID~v4 estándar
    (8-4-4-4-12 dígitos hexadecimales con el nibble de versión fijado a~4).
    Las peticiones con formato incorrecto reciben una respuesta \texttt{404 Not Found}.
  \item \textbf{Protección contra \textit{path traversal}}: el servicio verifica que el
    nombre del fichero solicitado no contenga secuencias como \texttt{../} que pudieran
    permitir el acceso a ficheros fuera del directorio de simulaciones configurado. Los
    intentos de \textit{path traversal} se rechazan con \texttt{404 Not Found} sin
    revelar información sobre la estructura del sistema de ficheros.
\end{itemize}

% ============================================================
\section{Diseño de la API REST}
\label{sec:api}
% ============================================================

La API REST constituye la interfaz de comunicación entre el backend y cualquier cliente que
desee consumir los resultados de las simulaciones. El diseño sigue los principios REST
clásicos: recursos identificados por URI, operaciones sin estado y representaciones en
JSON.

\subsection{Definición de Endpoints}
\label{subsec:endpoints}

El Cuadro~\ref{tab:endpoints} resume los \textit{endpoints} disponibles en la API. Todos
ellos son operaciones de solo lectura (\lstinline{GET}), coherente con el hecho de que las
simulaciones se generan en MATLAB y el backend actúa exclusivamente como capa de consulta.

\begin{table}[H]
\centering
\caption{\textit{Endpoints} de la API REST.}
\label{tab:endpoints}
\small
\begin{tabularx}{\textwidth}{|l|X|c|}
\hline
\textbf{Endpoint} & \textbf{Descripción} & \textbf{Resp.} \\
\hline
\texttt{GET /api/simulations} &
  Devuelve el listado de simulaciones disponibles, incluyendo \texttt{id},
  \texttt{name}, \texttt{createdAt} y \texttt{metadata}. No incluye los
  \texttt{coveragePoints} para minimizar el tamaño de la respuesta. Soporta
  paginación opcional con los parámetros \texttt{?page=0\&size=10}. & 200 \\
\hline
\texttt{GET /api/simulations/\{uuid\}} &
  Devuelve el detalle completo de una simulación, incluyendo los
  \texttt{coveragePoints} (hasta 3600 objetos). & 200 / 404 \\
\hline
\texttt{GET /api/simulations/\{uuid\}/export} &
  Descarga el fichero JSON original de la simulación como adjunto
  (\texttt{Content-Disposition: attachment}). & 200 / 404 \\
\hline
\texttt{GET /swagger-ui.html} &
  Documentación interactiva OpenAPI generada automáticamente por
  \texttt{springdoc-openapi}. & 200 \\
\hline
\end{tabularx}
\end{table}

\subsection{Códigos de Respuesta HTTP}
\label{subsec:http_codes}

La API utiliza un conjunto reducido de códigos de respuesta HTTP, suficiente para cubrir
los casos de uso de un servicio de solo lectura:

\begin{itemize}
  \item \textbf{200 OK}: la petición se ha procesado correctamente. El cuerpo contiene el
    recurso solicitado en formato JSON.
  \item \textbf{404 Not Found}: el UUID proporcionado no corresponde a ninguna simulación
    existente, tiene un formato inválido o el intento de acceso ha sido rechazado por la
    protección de \textit{path traversal}. El cuerpo contiene un mensaje de error
    genérico que no revela detalles internos.
  \item \textbf{500 Internal Server Error}: error inesperado del servidor (por ejemplo,
    fallo de lectura del sistema de ficheros). El manejador global de excepciones
    (\lstinline{@ControllerAdvice}) captura la excepción y devuelve un cuerpo JSON
    estructurado con el \textit{timestamp} y el mensaje de error.
\end{itemize}

\subsection{Documentación OpenAPI}
\label{subsec:openapi}

La documentación de la API se genera automáticamente mediante la biblioteca
\lstinline{springdoc-openapi} en su versión 2.3.0, declarada como dependencia Maven en el
fichero \lstinline{pom.xml} del proyecto. Esta biblioteca escanea las anotaciones de Spring
(\lstinline{@RestController}, \lstinline{@GetMapping}, \lstinline{@PathVariable}) y genera
una especificación OpenAPI~3.0 en tiempo de ejecución, accesible tanto en formato JSON
(\lstinline{/v3/api-docs}) como a través de la interfaz gráfica Swagger~UI en
\lstinline{/swagger-ui.html}.

La documentación interactiva permite al desarrollador o al integrador probar los
\textit{endpoints} directamente desde el navegador, visualizar los esquemas de respuesta y
consultar los códigos de error posibles. Esto elimina la necesidad de mantener una
documentación separada que pueda quedar desactualizada respecto al código fuente.

% ============================================================
\section{Diseño de la Interfaz de Usuario}
\label{sec:ui}
% ============================================================

La interfaz de usuario sigue una disposición de panel lateral con mapa principal,
optimizada para que el operador pueda seleccionar y comparar simulaciones sin perder de
vista el mapa de cobertura. El diseño se ha validado con \textit{mockups} de alta
fidelidad antes de la implementación.

\subsection{Componentes de la Aplicación}
\label{subsec:componentes}

La aplicación Angular se estructura en siete componentes \textit{standalone}, cada uno con
una responsabilidad claramente definida. El Cuadro~\ref{tab:componentes} los describe.

\begin{table}[H]
\centering
\caption{Componentes Angular de la aplicación.}
\label{tab:componentes}
\begin{tabular}{|l|p{8.5cm}|}
\hline
\textbf{Componente} & \textbf{Responsabilidad} \\
\hline
\texttt{MapComponent} &
  Inicializa el mapa Leaflet, gestiona las capas de \textit{tiles} OpenStreetMap y
  renderiza la capa de calor (\textit{heatmap}) mediante \texttt{leaflet.heat}. \\
\hline
\texttt{SidebarComponent} &
  Panel lateral con el selector de simulaciones, los controles de filtrado y el acceso
  a la configuración de visualización. \\
\hline
\texttt{HeaderComponent} &
  Barra superior con el nombre de la aplicación, el logotipo y el botón de alternancia
  de tema claro/oscuro. \\
\hline
\texttt{InfoPanelComponent} &
  Panel flotante que muestra los detalles del punto seleccionado al hacer clic sobre el
  mapa: coordenadas, valor RSRP y clasificación de cobertura. \\
\hline
\texttt{StatsPanelComponent} &
  Muestra estadísticas agregadas de la simulación activa: RSRP medio, mínimo y máximo,
  número de puntos por zona de cobertura y porcentaje de área cubierta. \\
\hline
\texttt{CoverageLegendComponent} &
  Leyenda visual con la escala de colores RSRP, indicando los umbrales y la
  clasificación de calidad de cada zona. \\
\hline
\texttt{SimulationConfigComponent} &
  Muestra los parámetros de la simulación seleccionada: frecuencia, potencia, modelo
  de propagación, fecha de ejecución y tipo de datos de terreno. \\
\hline
\end{tabular}
\end{table}

Además de los componentes visuales, la aplicación cuenta con dos servicios inyectables:
\lstinline{SimulationService}, que encapsula las llamadas HTTP al backend y mantiene el
estado reactivo mediante \textit{signals}, y \lstinline{ThemeService}, que gestiona la
alternancia entre tema claro y oscuro persistiendo la preferencia del usuario.

\subsection{Interacción con el Mapa}
\label{subsec:interaccion_mapa}

El mapa constituye el elemento central de la interfaz y ofrece tres formas de interacción:

\begin{enumerate}
  \item \textbf{Inspección por clic}: al hacer clic sobre cualquier punto del
    \textit{heatmap}, se activa el \lstinline{InfoPanelComponent}, que muestra las
    coordenadas geográficas del punto más cercano y su valor RSRP en dBm, junto con la
    clasificación de cobertura correspondiente.
  \item \textbf{Actualización automática}: cuando el usuario selecciona una simulación
    diferente en el \lstinline{SidebarComponent}, el \textit{heatmap} se recalcula y
    rerenderiza automáticamente. El componente de mapa destruye la capa de calor anterior,
    crea una nueva con los datos de la simulación seleccionada y ajusta la vista del mapa
    al \textit{bounding box} de los nuevos puntos.
  \item \textbf{Filtrado por umbral RSRP}: un control deslizante (\textit{slider}) permite
    al usuario establecer un umbral mínimo de RSRP. Los puntos con valores inferiores al
    umbral se ocultan del \textit{heatmap}, permitiendo visualizar únicamente las zonas con
    calidad de señal suficiente para un servicio concreto.
\end{enumerate}

\subsection{Paleta de Colores RSRP}
\label{subsec:paleta_colores}

La escala de colores utilizada para representar los niveles de RSRP en el \textit{heatmap}
se ha diseñado siguiendo las convenciones habituales en herramientas de planificación
radioeléctrica, con una progresión intuitiva de verde (señal fuerte) a rojo oscuro (sin
cobertura). Los umbrales se basan en las recomendaciones de 3GPP TS~38.215
\cite{3gpp_ts38215_2022} y en la práctica habitual de los operadores europeos.

\begin{table}[H]
\centering
\caption{Paleta de colores y umbrales RSRP.}
\label{tab:paleta_rsrp}
\footnotesize
\begin{tabularx}{\textwidth}{|c|c|l|X|}
\hline
\textbf{Umbral RSRP} & \textbf{Color} & \textbf{Clasificación} & \textbf{Descripción} \\
\hline
$\geq -80$~dBm     & Verde         & Excelente     & Señal óptima; todos los servicios
                                                      5G disponibles (eMBB, URLLC). \\
\hline
$-90$ a $-80$~dBm   & Amarillo-verde & Buena        & Calidad de servicio estable para
                                                      la mayoría de aplicaciones. \\
\hline
$-100$ a $-90$~dBm  & Naranja       & Aceptable     & Servicio básico de datos disponible;
                                                      posible degradación en servicios
                                                      de baja latencia. \\
\hline
$-110$ a $-100$~dBm & Rojo          & Débil         & Señal muy débil; conexión inestable
                                                      y alto riesgo de pérdida de enlace. \\
\hline
$< -110$~dBm        & Rojo oscuro   & Sin cobertura & Sin servicio 5G utilizable;
                                                      por debajo del umbral de
                                                      sensibilidad del receptor. \\
\hline
\end{tabularx}
\end{table}

% ============================================================
\section{Decisiones de Diseño}
\label{sec:decisiones}
% ============================================================

Esta sección documenta las decisiones técnicas más relevantes tomadas durante la fase de
diseño, junto con su justificación.

\subsection{Justificación del Stack Tecnológico}
\label{subsec:just_stack}

La elección de cada tecnología responde a criterios de adecuación funcional, madurez del
ecosistema y disponibilidad de herramientas de desarrollo:

\begin{itemize}
  \item \textbf{MATLAB} para la capa de simulación: MATLAB dispone del
    \textit{Communications Toolbox}, que implementa el modelo Longley-Rice ITM de forma
    nativa mediante \lstinline{propagationModel('longley-rice')}, junto con funciones de
    creación de emplazamientos (\lstinline{txsite}, \lstinline{rxsite}) y cálculo de
    potencia recibida (\lstinline{sigstrength}). El \textit{Mapping Toolbox} proporciona
    funciones de carga e interpolación de datos SRTM. Ningún otro entorno ofrece una
    integración equivalente entre modelado de propagación y datos de elevación con la misma
    productividad.

  \item \textbf{Spring Boot 3.2 + Java~17} para el backend: Spring Boot es el
    \textit{framework} Java más utilizado para servicios REST, con un ecosistema maduro de
    \textit{starters}, gestión automática de dependencias mediante Maven y generación
    de documentación OpenAPI con \lstinline{springdoc-openapi}. Java~17, al ser una versión
    LTS (\textit{Long-Term Support}), garantiza estabilidad y soporte a largo plazo. Los
    \textit{records} de Java~16+ permiten definir DTOs inmutables de forma concisa.

  \item \textbf{Angular 17+} para el frontend: Angular ofrece tipado estático mediante
    TypeScript, lo que reduce los errores en tiempo de ejecución y mejora la
    refactorización del código. Los componentes \textit{standalone} introducidos en
    Angular~14 y consolidados en Angular~17 simplifican la arquitectura del proyecto. El
    sistema de \textit{signals} proporciona reactividad eficiente sin dependencias externas
    de gestión de estado.

  \item \textbf{Leaflet + leaflet.heat} para la visualización cartográfica: Leaflet es la
    biblioteca de mapas web de código abierto más ligera y extendida, con un tamaño de
    apenas 42~kB. El \textit{plugin} \lstinline{leaflet.heat} añade capacidad de
    renderizado de \textit{heatmaps} con un API mínimo, ideal para representar los 3600
    puntos de cobertura RSRP sin sobrecarga de rendimiento.
\end{itemize}

\subsection{Selección del Modelo de Propagación}
\label{subsec:seleccion_modelo}

La selección del modelo de propagación es la decisión de diseño de mayor impacto técnico
en la precisión del gemelo digital. Se evaluaron tres modelos en un escenario de referencia
(campus URJC de Fuenlabrada, cuadrícula $60 \times 60$, frecuencia 3,5~GHz, EIRP 53~dBm).
Los resultados completos de la comparativa se presentan en el
Capítulo~\ref{cap:capitulo5}; aquí se resumen las conclusiones que fundamentan la decisión
de diseño.

\paragraph{FSPL (\textit{Free Space Path Loss}).} Se utiliza únicamente como referencia
teórica de cota superior. Al no incorporar terreno ni obstáculos, produce valores de
potencia recibida irrealmente elevados (hasta $-21,4$~dBm en zona cercana). No es un modelo
válido para planificación real, pero resulta útil como prueba de coherencia: cualquier
modelo realista debe producir pérdidas mayores que la FSPL.

\paragraph{COST-231 Hata.} Este modelo empírico está validado para el rango
1500--2000~MHz. Su extrapolación a 3,5~GHz introduce un error sistemático de
$+15$ a $+20$~dB en pérdidas, lo que produce resultados catastróficos: en el escenario de
referencia, el 95,2~\% del área aparece sin cobertura, un resultado incompatible con la
realidad de un emplazamiento con EIRP de 53~dBm a distancias inferiores a 1,5~km.
Se mantiene en el sistema exclusivamente como \textit{fallback} de emergencia para
entornos MATLAB sin \textit{Communications Toolbox}, con advertencia explícita en los
metadatos del JSON generado.

\paragraph{Longley-Rice ITM (modelo seleccionado).} El \textit{Irregular Terrain Model}
\cite{longley1968transmission}, validado por la NTIA para el rango 20~MHz--20~GHz
\cite{hufford1982guide}, es el único de los tres modelos que incorpora el perfil real de
elevación del terreno en el cálculo de pérdidas. En el escenario de referencia produce un
RSRP medio de $-93,9$~dBm, con un 93,4~\% de los puntos por encima del umbral de
cobertura aceptable ($-100$~dBm), lo que es coherente con las expectativas para una
macrocelda 5G suburbana. El Cuadro~\ref{tab:comp_modelos} resume la comparativa.

\begin{table}[H]
\centering
\caption{Comparativa de modelos de propagación (escenario URJC Fuenlabrada, 3,5~GHz).}
\label{tab:comp_modelos}
\begin{tabular}{|l|c|c|c|}
\hline
\textbf{Criterio} & \textbf{FSPL} & \textbf{COST-231} & \textbf{Longley-Rice} \\
\hline
Validez a 3,5~GHz       & Sí (teórico)   & No ($\pm$15~dB)     & \textbf{Sí} \\
Incorpora terreno        & No             & No                   & \textbf{Sí} \\
Coherencia con TR~38.901 & Referencia     & Parcial              & \textbf{Alta} \\
Disponibilidad MATLAB    & Siempre        & Siempre              & Comm.~Toolbox \\
RSRP medio               & ---            & $-125,8$~dBm         & $\mathbf{-93,9}$~\textbf{dBm} \\
Cobertura $\geq -100$~dBm & ---           & 1,4~\%               & \textbf{93,4~\%} \\
Rol en el sistema        & Validación     & \textit{Fallback}    & \textbf{Principal} \\
\hline
\end{tabular}
\end{table}

La decisión de diseño es, por tanto, utilizar el modelo Longley-Rice ITM como modelo
principal del gemelo digital, con degradación automática a COST-231 Hata ---debidamente
señalizada--- si el \textit{Communications Toolbox} no está disponible en el entorno de
ejecución.
